%% ============================================================
%%  chapters/02_psu_recruitment.tex
%%  Chapter 2: PSU Recruitment and Preparation
%% ============================================================

\chapter{PSU Recruitment \& Preparation}
\label{chap:psu}

\begin{protip}[Chapter Overview]
  Public Sector Undertakings (PSUs) offer stable, high-paying careers with strong
  growth prospects. This chapter covers the recruitment process, cutoffs, and
  strategic preparation tips sourced directly from PrepFusion students.
\end{protip}

%% ─────────────────────────────────────────────────────────
\section{Overview of PSU Recruitment Modes}
\label{sec:psu-overview}

PSU recruitment has evolved significantly in recent years. Two primary modes now
co-exist:

\begin{tblr}{
    colspec = {X[2,l] X[4,l] X[3,l]},
    hlines, vlines,
    row{1} = {bg=pfBlue, fg=white, font=\sffamily\bfseries},
    row{2,4} = {bg=pfGrey},
  }
  \toprule
  Mode & Description & Examples \\
  \midrule
  \textbf{\GATE Score Based}
    & Recruitment purely through valid \GATE rank, followed by GD \& interview.
    & IOCL, GAIL, NTPC, BHEL, BEL, PGCIL \\
  \textbf{Own Exam Based}
    & PSU conducts its own written examination followed by GD \& interview.
    & DRDO, some state PSUs \\
  \bottomrule
\end{tblr}

\vspace{6pt}

\begin{eligalert}[GATE Score Validity for PSUs]
  \begin{pfChecklist}
    \item \textbf{Most PSUs:} Accept current year \GATE rank only.
    \item \textbf{BARC:} Accepts \GATE scores up to \textbf{2 years} old.
    \item \textbf{NPCIL:} Accepts \GATE scores up to \textbf{3 years} old
          (final selection based entirely on technical interview).
  \end{pfChecklist}
\end{eligalert}

%% ─────────────────────────────────────────────────────────
\section{Recruitment Process (Step-by-Step)}
\label{sec:psu-process}

\begin{enumerate}[label=\textcolor{pfBlue}{\bfseries Stage \arabic*:}, itemsep=8pt, leftmargin=*]

  \item \textbf{Written Test / \GATE Shortlisting.}
        Candidates are shortlisted based on their \GATE rank or the PSU's own
        written exam score. For GATE-based PSUs, typically only the
        \textbf{current year} rank is valid.

  \item \textbf{Group Discussion (GD).}
        Many PSUs include a GD round to evaluate communication skills,
        situational awareness, and team dynamics.

  \item \textbf{Technical Interview.}
        Subjects tested vary by discipline. For EE/ECE:
        \begin{pfChecklist}
          \item Power Systems (protection, stability)
          \item Electrical Machines (synchronous, induction, transformers)
          \item Analog \& Digital Circuits
          \item Company-specific topics (business, chairman/MD name, projects)
        \end{pfChecklist}

  \item \textbf{HR / Personal Interview.}
        Expect questions on: Why PSU over private sector? Where do you see yourself
        in 10 years? Why this specific PSU?

  \item \textbf{Medical Examination.}
        A rigorous medical clearance is mandatory before final selection.
        Multiple tests must be cleared; failure at this stage results in
        disqualification.

\end{enumerate}

%% ─────────────────────────────────────────────────────────
\section{Score Weightage Breakdown}
\label{sec:psu-weightage}

\begin{tblr}{
    colspec = {X[3,l] X[2,c] X[2,c] X[2,c]},
    hlines, vlines,
    row{1} = {bg=pfBlue, fg=white, font=\sffamily\bfseries},
    row{2,4,6} = {bg=pfGrey},
  }
  \toprule
  PSU & \GATE Score & GD & Interview \\
  \midrule
  IOCL   & 80--85\% & \SetCell[c=2]{c} 15--20\% combined & \\
  GAIL   & 80--85\% & \SetCell[c=2]{c} 15--20\% combined & \\
  NTPC   & Rank-based shortlist & \SetCell[c=2]{c} Interview-heavy & \\
  BARC   & Qualifying only & \SetCell[c=2]{c} 100\% technical interview & \\
  NPCIL  & Qualifying only & \SetCell[c=2]{c} 100\% technical interview & \\
  \bottomrule
\end{tblr}

%% ─────────────────────────────────────────────────────────
\section{Rank-wise PSU Expectation (ECE, OBC Category)}
\label{sec:psu-rank}

\begin{tblr}{
    colspec = {X[2,l] X[3,l] X[4,l]},
    hlines, vlines,
    row{1} = {bg=pfBlue, fg=white, font=\sffamily\bfseries},
    row{2,4,6} = {bg=pfGrey},
  }
  \toprule
  \GATE Rank Range & Targetable PSUs & Notes \\
  \midrule
  Under 150 (General)
    & Top PSUs (IOCL, GAIL, BHEL, PGCIL, NTPC)
    & Safe zone for general category in most PSUs. \\
  150--500
    & NTPC, BEL, BHEL, CIL, NPCIL
    & Apply broadly; selection depends on GD/interview. \\
  1300--1400 (OBC)
    & NTPC (Round 2, cutoff $\approx$1118), CIL (Round 3)
    & NPCIL 2024 OBC cutoff was $\approx$563 \GATE score. \\
  40 marks ($\approx$560 score)
    & Selected NITs via CCMT; BITS Pilani (HD program); PhD interviews at IITs
    & Visit CCMT website for NIT VLSI branch cutoffs. \\
  \bottomrule
\end{tblr}

%% ─────────────────────────────────────────────────────────
\section{Career Growth in PSUs}
\label{sec:psu-career}

\begin{tblr}{
    colspec = {X[2,l] X[5,l]},
    hlines, vlines,
    row{1} = {bg=pfBlue, fg=white, font=\sffamily\bfseries},
    row{2,4,6} = {bg=pfGrey},
  }
  \toprule
  Years of Service & Typical Designation \\
  \midrule
  0--5 years   & Executive / Junior Manager \\
  5--15 years  & Senior Manager / AGM \\
  15--25 years & DGM / GM \\
  25--30+ years & ED / Director / Chairman (rare, high-performers) \\
  \bottomrule
\end{tblr}

\begin{protip}[PSU Career Advantages]
  Beyond promotions, many PSUs:
  \begin{pfChecklist}
    \item Assign responsibilities \textit{outside} your engineering domain,
          broadening your skill set.
    \item Sponsor training sessions run by top technology companies, yielding
          industry-recognized certifications.
    \item Offer excellent work-life balance, medical benefits, and housing
          allowances.
  \end{pfChecklist}
\end{protip}

%% ─────────────────────────────────────────────────────────
\section{PSU Interview Preparation Strategy}
\label{sec:psu-prep}

\subsection{Technical Preparation}

\begin{pfChecklist}
  \item \textbf{Power Systems:} Protection schemes, stability analysis (for EE).
  \item \textbf{Electrical Machines:} Synchronous machines, induction motors,
        transformers---ratings, nameplate data, working principles.
  \item \textbf{Analog \& Digital:} Core ECE topics; expect basic definition
        questions (e.g., ``Define a Flip-Flop'').
  \item \textbf{Final Year Project:} Especially for PSUs like \textbf{Powergrid}
        and \textbf{EIL}, which deep-dive into your project. Prepare your BIODATA
        accordingly.
  \item \textbf{BEL/BHEL Strategy:} Practice previous-year questions (easier
        than GATE). Solve in a time-bound manner. Focus on practical aspects and
        theory.
\end{pfChecklist}

\subsection{Non-Technical / HR Preparation}

\begin{pfChecklist}
  \item Research the company: Chairman, MD, major projects, financial figures.
  \item Prepare answers for: ``Why PSU over a private company?''
  \item Prepare answers for: ``Why not M.Tech if you are a fresher?''
  \item Research papers and internships are \textbf{less critical} for PSU
        interviews, except for EIL where practical experience matters.
\end{pfChecklist}

\subsection{M.Tech vs.\ Direct PSU}

\begin{eligalert}[M.Tech On-Campus Placement Rule]
  M.Tech on-campus placement is based \textbf{entirely on M.Tech CGPA}.
  Your B.Tech CGPA has \textbf{no role} in campus recruitment during M.Tech.
  This is a significant advantage for candidates with a lower B.Tech GPA.
\end{eligalert}

%% ─────────────────────────────────────────────────────────
\section{Frequently Asked Questions}
\label{sec:psu-faq}

\begin{faqentry}{What is the weightage of research papers and final year
  projects in PSU interviews?}
  Research papers are generally \textbf{not important} for most PSU interviews.
  However, practical industry experience, internships, and final year projects
  play a significant role in some PSUs like \textbf{EIL}. For Powergrid,
  interviewers deep-dive into the projects mentioned on your BIODATA form.
\end{faqentry}

\begin{faqentry}{What are the opportunities for a candidate scoring 40 marks
  in GATE ECE 2025?}
  In GATE 2024, 40 marks corresponded to approximately a 560 score. At this
  level, PhD interview opportunities exist at IITs/IISc (interview cutoffs are
  very low). CCMT provides access to some NIT VLSI-related branches. BITS Pilani
  HD programs are also a good option if fees are not a constraint.
\end{faqentry}
